\documentclass[article,nojss]{jss}

%%%%%%%%%%%%%%%%%%%%%%%%%%%%%%
%% declarations for jss.cls %%%%%%%%%%%%%%%%%%%%%%%%%%%%%%%%%%%%%%%%%%
%%%%%%%%%%%%%%%%%%%%%%%%%%%%%%

%% almost as usual
\author{M\r{a}ns Magnusson\\Aalto University \And
        Markus Kainu\\Plus Affiliation \And
        Janne Huovari\\Plus Affiliation \And
        Leo Lahti\\University of Turku}

%\title{\pkg{pxweb}: a toolkit for open public data}
\title{Opening Up Official Statistics with the \pkg{pxweb} Package}

%% for pretty printing and a nice hypersummary also set:
\Plainauthor{Mans Magnusson, Markus Kainu, Janne Huovari, Leo Lahti} %% comma-separated
\Plaintitle{Opening Up Official Statistics with the pxweb Package} %% without formatting
\Shorttitle{\pkg{pxweb}: Opening Up Official Statistics} %% a short title (if necessary)

%% an abstract and keywords
\Abstract{
  The \pkg{pxweb} package provides an interface from
  \proglang{R} to statistical databases published through the PX-Web
  API. PX-Web is used by national, regional, and international
  statistical organizations to disseminate official statistics in
  multidimensional tables. The package supports discovery of API
  resources, construction and validation of table queries, retrieval and
  parsing of data and metadata, conversion to standard \proglang{R}
  data structures, extraction of comments and codelists, and citation of
  retrieved data. This article introduces the motivation and design of
  \pkg{pxweb}, outlines the support for both PX-Web API version 1 and
  version 2, and provides an reproducible case study using official statistics.
}
\Keywords{open data, citing data, public data, pxweb, PC-Axis, \proglang{R}}
\Plainkeywords{open data, citing data, public data, pxweb, PC-Axis, R} %% without formatting
%% at least one keyword must be supplied

%% Compatibility aliases while the manuscript is migrated from the R Journal
%% template to the JSS template.
\newcommand{\CRANpkg}[1]{\pkg{#1}}
\newcommand{\Cpkg}[1]{\pkg{#1}}

%% publication information
%% NOTE: Typically, this can be left commented and will be filled out by the technical editor
%% \Volume{50}
%% \Issue{9}
%% \Month{June}
%% \Year{2012}
%% \Submitdate{2012-06-04}
%% \Acceptdate{2012-06-04}

%% The address of (at least) one author should be given
%% in the following format:
\Address{
  M\r{a}ns Magnusson\\
  Department of Statistics\\
  Uppsala University\\
  Sweden\\
  E-mail: \email{mons.magnusson@gmail.com}\\
\\
  Leo Lahti\\
  Department of Mathematics and Statistics\\
  University of Turku\\
  Finland\\
  E-mail: \email{leo.lahti@iki.fi}
}

%% It is also possible to add a telephone and fax number
%% before the e-mail in the following format:
%% Telephone: +43/512/507-7103
%% Fax: +43/512/507-2851

%% for those who use Sweave please include the following line (with % symbols):
%% need no \usepackage{Sweave.sty}

%% end of declarations %%%%%%%%%%%%%%%%%%%%%%%%%%%%%%%%%%%%%%%%%%%%%%%


\begin{document}

%% This is now the canonical JSS manuscript source. Earlier R Journal
%% draft material is retained in magnusson-kainu-lahti.tex for reference.


\section{Introduction}

Open data science workflows rely heavily on algorithmic tools for data
retrieval. Automating major parts of the data workflow,
such as finding, accessing, integrating, citing, and reporting data,
helps users spend more time on statistical analysis and interpretation.
It can also improve reproducible research \citep{Gandrud13,
Boettiger2015} by making data acquisition explicit and repeatable.

Official statistics are a particularly important class of public data.
National statistical institutes and international agencies publish
large collections of data on demography, economics, health,
infrastructure, climate, and other areas \citep{TODO,TODO}. Such data are often
multidimensional, geographically detailed, and updated over long time
periods. Opening these resources through public web interfaces is,
however, only the first step. Users also need software that can bridge
the gap between online statistical databases and reproducible analysis
in statistical programming environments such as \proglang{R}.

Dedicated software packages simplify, standardize, and automate
analysis workflows. They take into account variations in raw data
formats, access details, metadata conventions, and typical use cases so
that users can avoid repetitive programming tasks and reduce the risk of
misinterpretation and data processing errors \citep{TODO}.
In \proglang{R}, several packages support generic or
source-specific data retrieval, including \pkg{quandl} \citep{quandl},
\pkg{pdfetch} \citep{pdfetch}, \pkg{eurostat}
\citep{Lahti17eurostat}, \pkg{WDI} \citep{WDI}, and \pkg{osmar}
\citep{osmar}.

The \pkg{pxweb} package fills a complementary role: it provides a
general \proglang{R} interface to databases exposed through the PX-Web
API. The package can navigate PX-Web resources, fetch metadata,
construct and validate queries, retrieve data, convert responses into
standard \proglang{R} data structures, and produce citation information
for the source data. As of version 1.0.0, the
package includes 46 bundled API catalogue entries and supports both
PX-Web API version 1 and version 2.

Data citation is another motivation for the package. Data reuse should
document source organizations, accessed resources, version information,
and retrieval times where possible. \pkg{pxweb} helps standardize this
process for PX-Web data by collecting source information from the API
catalogue and data objects and by encouraging users to cite both the
retrieved data and the software used to retrieve them.

The package was first released on CRAN in 2014 and was substantially
reworked in 2018 and 2026. Recent development has added support for the PX-Web
API version 2, including table search, JSON-stat2 responses, v2
metadata, codelists, value sets, aggregation helpers, and observation
status comments. The current package is part of the rOpenGov open data
science project \citep{Lahti13icml}.

\section[PX-Web and PC-Axis]{PX-Web and PC-Axis}
\label{sec:pxweb-pcaxis}

PX-Web is a web application and API for disseminating PC-Axis
statistical data \citep{TODO}. PC-Axis data are typically organized as
multidimensional tables with dimensions such as region, time, sex, age,
indicator, or content variable and has a long tradition in national
statistical agencies \citep[NSI, ][]{TODO}. PX-Web exposes these data and their
metadata through HTTP endpoints that can be explored interactively or
accessed programmatically.

From the perspective of an \proglang{R} user, PX-Web APIs have three
important purposes. First, a table query must be constructed from table
metadata. Second, the data response has to be parsed back into an
analysis-ready format while preserving enough metadata to make the
result interpretable. Third, the downloaded data should be easy to
cite to facilitate reproducibility. \pkg{pxweb} handles all three.

The package distinguishes between the PX-Web API and the older PC-Axis
file format, typically stored as \code{.px} files \citep{TODO}. \pkg{pxweb} is
designed for APIs, while packages such as \pkg{pxR} focus on parsing
legacy PC-Axis files \citep{TODO}. This distinction is important because API access
requires metadata discovery, query construction, request validation,
and response parsing rather than only local file import.

% TODO: Add a concise historical paragraph on PX-Web/PC-Axis and cite
% authoritative documentation from Statistics Sweden or the PX-Web project.
% TODO: Clarify current PX-Web API v1/v2 adoption among catalogue entries.
% TODO: Describe provider-side query limits and rate limits using current docs.

\section{Package Design}
\label{sec:package-design}

The \pkg{pxweb} package is designed as a generic interface rather than a
client for one statistical agency. Its core design goal is to provide a
stable \proglang{R} workflow across PX-Web providers while preserving
provider-specific metadata where it affects query construction,
codelists, comments, and citation.

The current package interface has 39 exported functions. The central
workflow is built around \code{pxweb_get()}, which retrieves hierarchy
levels, table metadata, or table data depending on the supplied URL and
query. The helper \code{pxweb_get_data()} wraps the same workflow and
returns a \code{data.frame}. Query construction is handled by
\code{pxweb_query()} and related conversion helpers, including
\code{pxweb_query_as_json()} and \code{pxweb_query_as_rcode()}.

Internally, the package represents API URLs, metadata, data, comments,
and query objects with S3 classes. This allows common operations such as
printing, validation, conversion with \code{as.data.frame()}, and
combining split responses to behave consistently across API versions.
The package uses \pkg{httr} for HTTP requests and \pkg{jsonlite} for
JSON parsing \citep{httr,jsonlite}.

The bundled API catalogue is a practical layer on top of the generic
API client. It stores provider URL templates, available languages,
version paths, aliases, and citation metadata. Users can rely on the
catalogue for common providers or pass explicit PX-Web URLs for APIs
that are not bundled with the package. The current catalogue contains
46 entries covering national, regional, and international statistical
resources, see Table \ref{TODO}.

Support for PX-Web API version 2 is expressed in the same top-level
workflow as version 1 \citep{TODO}. Version 2 tables use separate metadata and data
endpoints and return JSON-stat2 by default. The \pkg{pxweb} package detects v2 URLs,
builds the corresponding data endpoint, constructs the v2 query body,
and adds query parameters for codelists and output values when needed.
This keeps user code close to the v1 workflow while still exposing v2
features.

Large queries can exceed provider-side limits. Here, \pkg{pxweb} validates
queries against metadata and splits supported large requests into
batches when possible. Returned pieces are then combined into a single
data object where the response type supports it. This behavior is
important for official statistics tables with many regions, time
periods, or classifications that often can be larger than the actual
provider-side limits.

The package follows common open-source software practices, including
version control, automated tests, mocked API fixtures, continuous
integration, and live smoke-test workflows for selected APIs
\citep{PerezRiverol2016}. The default local test suite currently uses
mocked and offline-safe tests, with live tests controlled by
environment variables.

% TODO: Add a figure showing the object/workflow model: catalogue or URL,
% metadata, query, data request, parsed data, comments, citation.
% TODO: Add a compact table of exported functions grouped by workflow.

\section{Functionality}
\label{sec:functionality}

The package supports both interactive and production workflows.
Interactive use starts with \code{pxweb_interactive()}, which allows a
user to navigate a known API, select a table, build a query, optionally
download data, and print reusable \proglang{R} code. Code use
starts from either a catalogue entry or a URL and combines metadata
inspection, query construction, validation, and retrieval.

Installation of the CRAN release version follows the standard procedure in R.

\begin{Code}
install.packages("pxweb")
library("pxweb")
\end{Code}

In PX-WEB APIs version 2, users can search for tables with
\code{pxweb_search()}, retrieve metadata with \code{pxweb_get()}, build
a query as a named list, fetch data with \code{pxweb_get()} or
\code{pxweb_get_data()}, and convert the result to a \code{data.frame}.

\begin{Code}
url <- "https://statistikdatabasen.scb.se/api/v2/tables/TAB638/metadata?lang=sv"

query <- list(
  Region = "00",
  Civilstand = "OG",
  Alder = "0",
  Kon = "1",
  ContentsCode = "BE0101N1",
  Tid = "2024"
)

px_data <- pxweb_get(url, query = query)
px_df <- as.data.frame(
  px_data,
  column.name.type = "code",
  variable.value.type = "code"
)
\end{Code}

Version 2 APIs may expose aggregation codelists and value sets. The
package provides query helpers for all values, latest values, top and
bottom selections, aggregations, and value sets so that users can
express these selections in \proglang{R} rather than manually
constructing API-specific URL parameters. Available codelists can be
inspected with \code{pxweb_codelists()}.

\begin{Code}
meta <- pxweb_get(url)
pxweb_codelists(meta, variable = "Alder")

query <- list(
  Region = pxweb_all(),
  Alder = pxweb_aggregation("agg_Ålder5år_1"),
  Kon = c("1", "2"),
  ContentsCode = "BE0101N1",
  Tid = pxweb_latest()
)
\end{Code}

PX-Web responses may include comments, notes, or status information.
\pkg{pxweb} exposes these through \code{pxweb_data_comments()}, with
support for both v1 comments and v2 dataset notes, dimension notes,
value notes, and observation status comments. Citation support is
provided by \code{pxweb_cite()}, which prints citation information for
the retrieved data source and the \pkg{pxweb} package itself.

% TODO: Turn this section into a tighter software tour with one running
% example and short subsections for search, query, data conversion,
% comments/codelists, citation, and extension to uncatalogued APIs.

\section{Case Study}
\label{sec:case-study}

% TODO: Select one case study and implement it as JSS replication code.
% TODO: Add one generated table or figure and reference it from the text.

\section{Related Software}
\label{sec:related}

\pkg{pxweb} sits between general data-retrieval packages and
source-specific official-statistics packages. Packages such as
\pkg{quandl} and \pkg{pdfetch} provide access to broad data services or
time-series sources \citep{quandl,pdfetch}. Packages such as
\pkg{eurostat} and \pkg{WDI} focus on specific data providers
\citep{Lahti17eurostat,WDI}. In contrast, \pkg{pxweb} targets a shared
API family used by multiple statistical organizations.

The package should also be distinguished from tools for local PC-Axis
files. \pkg{pxweb} targets web APIs, including metadata discovery,
query validation, request splitting, and response parsing. File-oriented
tools such as \pkg{pxR} address a related but different problem: reading
local \code{.px} files.

Finally, \pkg{pxweb} can serve as infrastructure for more specialized
packages and applications. Domain-specific tools can build on its API
access, query, and parsing functionality while adding provider-specific
analysis, visualization, or educational workflows.

% TODO: Add current references for PxWebApiData, geofi, pxR, and any
% official PX-Web clients or agency-specific R packages.

\section{Discussion}
\label{sec:discussion}

The \pkg{pxweb} package provides access to official
statistics shared through the PX-Web API. It helps bridge the gap
between data providers and users by supporting discovery, query
construction, data retrieval, parsing, comments, codelists, and
citation. This enables transparent and reproducible workflows from raw
online data to statistical summaries and publication-ready outputs.

The package has matured through active development since its first CRAN
release in 2014, including a major redesign in 2018 and 2026, with recent support
for PX-Web API version 2. The current implementation is tested with a
combination of unit tests, mocked API responses, and optional live
smoke tests. The local default test suite currently passes 667 tests
with live tests skipped unless explicitly enabled.

Limitations remain. PX-Web installations differ across providers,
provider-side query limits vary, and live APIs can change independently
of the package. For this reason, robust metadata validation, mocked
fixtures, and live smoke tests are all important parts of the
maintenance strategy. The package also needs careful documentation of
which v2 features are generic and which depend on provider-specific
metadata conventions.

The package has been used, for instance, in research on electronic
gambling machines and socioeconomic status \citep{Raisamo2019}. More
generally, programmatic access to official statistics can support open
data science, governmental analysis, research, education, and citizen
science. The work contributes to the broader open data science
ecosystem \citep{Lahti2018IDA} by making official statistics easier to
retrieve, document, and reuse in \proglang{R}.

% TODO: Update with more citations and where the package has been used..
% TODO: Update download counts and usage indicators from current CRAN logs.
% TODO: Add a concise limitations paragraph based on current open issues.
% TODO: Replace this discussion with a stronger conclusion after the case
% study and related-software sections are finalized.

\section*{Acknowledgments}

We are grateful to all package contributors. The work has been
partially funded by Academy of Finland (decisions 295741, 345630 to
LL), and is part of rOpenGov\footnote{\url{https://github.ropengov.io}}.


\bibliography{magnusson-kainu-lahti}

\end{document}










